\section{Conclusion}
\label{sec:conclusion}

Under one preregistered protocol, with matched inputs, parameter counts, budgets and tuning,
a diffusion policy over action chunks did not outperform a deterministic MLP that sees the
same observation window and predicts the same chunks, on either Lift or Can. Its advantage
over a one-step MLP on Can is not separated from that of an LSTM or of the chunked MLP, so
the data give no evidence for the generative model beyond temporal context on these tasks,
and even the separation from the one-step MLP depends on how the uncertainty over three seeds
is computed. Diffusion was the most robust family to a large joint offset on Lift; on Can its
drop was not separated from the chunked MLP's, and part of its robustness may come from its
sampler's bound on its actions. It cost more than an order of
magnitude more training and inference time than the chunked MLP. With images, visual BC and
visual diffusion both reach the Lift ceiling, and our DP3-inspired point-cloud policy, which
uses world-frame coordinates and no augmentation, generalized clearly worse than the matched
RGB policy to unseen cube positions.

The most useful next steps follow from the limitations: more training seeds, harder tasks
that are not saturated, more than one rollout per start, a BC baseline with the same action
bound as the diffusion sampler, and a point-cloud arm with augmentation or an object-centred
frame. Each would need its own preregistration and freeze before it is run.
